% !TeX root = ../../main.tex
\subsection{Comparing the kappa and Hodge backgrounds}
\label{subsec:s08:05}

The relation $12\lambda_1=\kappa_1+\delta$ changes the entire
background, not just one coefficient. We first estimate a single
boundary insertion, summing over every genus split. A monotone,
log-concave sequence of mixed intersections will then control all
boundary insertions at once.

For $g\geq2$ and $0\leq h\leq g$, set
\begin{equation}\label{eq:s08:node-definition}
 T_{g,h}^{\delta}
   =\int_{\MM_g}\lambda_{g-h}\delta e^{\kappa_1}.
\end{equation}
As usual the top-degree part is selected by the integral.

\begin{lmm}[A single node]\label{lmm:s08:one-node}
For these indices one has the exact decomposition
\begin{equation}\label{eq:s08:one-node}
 \begin{aligned}
 T_{g,h}^{\delta}
 ={}&
 \begin{cases}
  \tfrac12 T_{g-1,h-1,2},&h\geq1,\\
  0,&h=0
 \end{cases}\\
 &+\frac12\sum_{G=1}^{g-1}
   \ \sum_{\substack{0\leq i\leq h\\ i\leq G,\ h-i\leq g-G}}
                    T_{G,i,1}\,T_{g-G,h-i,1}.
 \end{aligned}
\end{equation}
For every fixed $h\geq0$, as $g\to\infty$,
\begin{equation}\label{eq:s08:node-bound}
 |T_{g,h}^{\delta}|=O_h(g^{-1}T_{g,h}).
\end{equation}
\end{lmm}
\begin{prf}
\emph{The restriction formula.}
On the normalization of a one-node stratum, the logarithmic
dualizing line restricts to the logarithmic dualizing line of each
component, with the branches counted as markings. Pushing its
square down the universal curve shows that $\kappa_1$ restricts
additively. Equation~\eqref{eq:s02:hodge-factorization} gives the
corresponding Hodge restriction.

For a nonseparating node the normalization has genus $g-1$ and two
markings. Its class $\lambda_{g-h}$ has defect $h-1$ and vanishes
when $h=0$. For a separating node the two defects sum to $h$,
giving the inner range in \eqref{eq:s08:one-node}. An unmarked
one-valent component cannot be rational, so $1\leq G\leq g-1$.
The first factor $1/2$ exchanges the two branches of the loop.
The second accounts for ordering the two separating components;
it also gives the correct factor when their genera coincide.
These are the ordered-branch stack conventions of
\eqref{eq:s02:ordered-boundary}. Thus no additional factor is
hidden in the vertex integrals.

\smallskip\noindent
\emph{The full genus sum.}
For this estimate abbreviate $D=D_{g,h}$. The degree of the
nonseparating vertex is $D-1$. Its bound in
\eqref{eq:s08:marked-bound}, the positive asymptotic
\eqref{eq:s08:intersection-asymptotic}, and the bounded ratio
$\beta_{g-1}/\beta_g$ give $O_h(T_{g,h}/D)$.

In a separating term put
$a=2G-2+i$ and $c=2(g-G)-2+h-i$. They are nonnegative integers
and $a+c=D-1$. The uniform comparison for $\beta_G$ in
Theorem~\ref{thm:s06:lambda-g} gives
$\beta_G\beta_{g-G}\leq C\beta_g$, independently of $G$.
Since $(a+1)(c+1)\geq D$, the two marked bounds yield
\begin{equation}\label{eq:s08:node-convolution}
 \frac{|T_{G,i,1}T_{g-G,h-i,1}|}{T_{g,h}}
 \leq \frac{C_h}{D\binom{D-1}{a}}.
\end{equation}
For every integer $N\geq0$,
$\sum_{a=0}^N\binom Na^{-1}\leq4$. For $N\geq4$, separate
the two endpoints and their neighbors; every remaining binomial
coefficient is at least $N(N-1)/2$, which proves the bound.
The smaller $N$ are immediate. For each fixed $i$, the values of
$a$ coming from the allowed genera form a subset of
$\{0,\ldots,D-1\}$. There are at most $h+1$ choices of $i$.
Summing \eqref{eq:s08:node-convolution} therefore proves
\eqref{eq:s08:node-bound} for the whole boundary, including all
small components.
\end{prf}

We next use nefness of the full Hodge bundle. A vector bundle is
nef here when its tautological quotient line bundle is nef.
The needed semipositivity applies on the compactification, including
test curves contained in the boundary.

\begin{lmm}[Nef classes on the compactification]\label{lmm:s08:nef}
The Hodge bundle $\EE$ and its determinant are nef on $\MM_g$,
and $\kappa_1=12\lambda_1-\delta$ is ample as a rational divisor
class. Intersection inequalities for nef divisors on smooth
projective varieties may be applied after a common smooth
projective generically finite cover of $\MM_g$ and its Hodge
projective bundle.
\end{lmm}
\begin{prf}
Fujita's semipositivity theorem \cite[Main Theorem (0.6)]{Fuj78}
applies to the relative dualizing direct image over a smooth
projective curve. The stable family is replaced by a semistable
resolution; its Hodge bundle is unchanged. Kawamata's
\cite[Theorem 5]{Kaw81} gives the corresponding statement over a
smooth projective base with a normal-crossing discriminant and
unipotent local monodromy. These hypotheses hold after semistable
reduction; quotient-line nonnegativity descends under finite
covers of a test curve.

For a test curve contained generically in the boundary, first make
a finite cover labeling its generic graph and lift it to the
product of the component moduli stacks. The component families
have smooth generic fibers. Their Hodge bundles are nef by the
preceding result. Partial normalization gives the exact sequence
\eqref{eq:s02:hodge-normalization}, with a trivial cycle-space
quotient on this labeled cover. Extensions and direct sums of nef
bundles on a curve are nef, so the original Hodge bundle is nef as
well. This is the boundary reduction in
\cite[Proposition 4.4]{LX26b}. Its determinant is nef: on each
test curve its degree is the top self-intersection of the
tautological nef class on the projective bundle.

The ampleness assertion follows from
\cite[Theorem (1.3)]{CH88}: for real coefficients $a,c$,
$a\lambda_1-c\delta$ is ample precisely when $a>11c>0$.
Take $a=12$, $c=1$, and use
\eqref{eq:s02:lambda-one-grr} with no markings.
Finally \cite[Proposition 6.1, $m=0$]{LX26b} supplies a smooth
projective generically finite cover of the stack. Pulling back
$\EE$ and then taking its projective bundle gives the corresponding
smooth projective cover upstairs. Top intersections are multiplied
by their common positive degree, including the generic stack
stabilizer factor, so inequalities descend by division. Pullback
preserves nefness; ampleness after this generically finite
pullback is not needed.
\end{prf}

For the remainder of this subsection let
$p:\mathbb P(\EE)\longrightarrow\MM_g$ parametrize one-dimensional
quotients, and put $\xi=c_1(\mathcal O_{\mathbb P(\EE)}(1))$.
Every base class in an intersection on this projective bundle is
pulled back by $p$.

\begin{prp}[Replacing the background]\label{prp:s08:background}
For each fixed integer $h\geq0$, as $g\to\infty$,
\begin{equation}\label{eq:s08:background}
 b_g(h)=12^{-D_{g,h}}D_{g,h}!\,T_{g,h}
                                      (1+O_h(g^{-1})).
\end{equation}
\end{prp}
\begin{prf}
\emph{Chern classes as projective intersections.}
For a local formal variable $t$, the projective bundle formula and
\eqref{eq:s02:chern-exponential} give
\begin{equation}\label{eq:s08:projective-pushforward}
 p_*\xi^{g-1+j}
 =[t^j]\,c_{-t}(\EE)^{-1}
 =[t^j]\,c_t(\EE)=\lambda_j
                                      \qquad(0\leq j\leq g).
\end{equation}
The middle equality uses the absence of even Chern characters:
$c_t(\EE)c_{-t}(\EE)=1$. It is a special Hodge-bundle identity,
not an identification of arbitrary Segre and Chern classes.

In this proof put $D=D_{g,h}$. For integers $0\leq k\leq D$ define
the mixed intersections
\begin{equation}\label{eq:s08:mixed-intersections}
 C_k=\int_{\mathbb P(\EE)}
       \xi^{2g-h-1}\,p^*(12\lambda_1)^k\,p^*\kappa_1^{D-k}.
\end{equation}
Their total degree is $4g-4=\dim\mathbb P(\EE)$. All three divisor
classes are nef by Lemma~\ref{lmm:s08:nef}. Since $\delta$ is
effective, their intersections on its pullback show that
$C_{k+1}\geq C_k$. The pushforward formula identifies their
endpoints and first difference:
\begin{equation}\label{eq:s08:mixed-endpoints}
 C_0=D!\,T_{g,h},\qquad
 C_D=12^D b_g(h),\qquad
 C_1-C_0=(D-1)!\,T_{g,h}^{\delta}.
\end{equation}
Here $D\geq1$ for every unmarked stable target under consideration.
The positive asymptotic in
\eqref{eq:s08:intersection-asymptotic} ensures $C_0>0$ for all
sufficiently large $g$.

\smallskip\noindent
\emph{The mixed sequence is log-concave.}
On the smooth projective cover from Lemma~\ref{lmm:s08:nef},
approximate each of the three nef classes by adding a small
positive multiple of a fixed ample class. For rational
approximations, take sufficiently divisible very ample multiples.
For $1\leq k<D$, cut a smooth surface by
$2g-h-1$ copies of the approximated $\xi$,
$k-1$ copies of the approximated $p^*(12\lambda_1)$, and
$D-k-1$ copies of the approximated $p^*\kappa_1$.
The Hodge index theorem on that surface gives
$C_k^2\geq C_{k-1}C_{k+1}$ after the approximation tends to zero.
Dividing by the common cover degree preserves this homogeneous
inequality. Thus the positive ratios $C_{k+1}/C_k$ decrease.

\smallskip\noindent
\emph{All boundary powers at once.}
Using monotonicity, log-concavity and
\eqref{eq:s08:mixed-endpoints}, we obtain
\begin{equation}\label{eq:s08:background-squeeze}
 1\leq\frac{C_D}{C_0}
 \leq\left(\frac{C_1}{C_0}\right)^D
 =\left(1+\frac{T_{g,h}^{\delta}}{D\,T_{g,h}}\right)^D
 \leq \exp\left(\frac{T_{g,h}^{\delta}}{T_{g,h}}\right)
 =1+O_h(g^{-1}).
\end{equation}
Lemma~\ref{lmm:s08:one-node} gives the last estimate.
The endpoint identities \eqref{eq:s08:mixed-endpoints} therefore
prove \eqref{eq:s08:background} for the full background.
\end{prf}
